\section{模型假设与符号说明}

\subsection{模型假设}

结合题意和第 2 节的分析，本文作如下假设。

（1）地球视为质点，卫星只受地球中心引力作用，轨道运动服从二体模型，不计 $J_2$ 项、日月引力和太阳光压等摄动。题目明确给定了这一动力学模型，30 天内摄动引起的轨道变化对观测几何的影响也较小。

（2）轨道机动为瞬时脉冲，机动前后位置连续、速度发生跳变，不考虑推力器的执行误差、指向误差和卫星姿态约束。同步轨道上一次机动的燃烧时间通常远小于轨道周期，用脉冲近似是合理的。

（3）太阳位置取 IAU 标准天文库 SOFA 给出的地心几何位置，只用于计算光照角，不产生引力作用。光照条件只按题目给定的光照角约束判断，不另行考虑地影。

（4）只要观测几何在连续 60~s 内满足题目的三个条件，即认为该次观测有效，不再考虑载荷视场、成像曝光等其他限制。同一目标被多次观测时只计一次。

（5）“每次有效观测最多 2 次脉冲”按全局时间轴理解，即从上一次观测结束到本次观测开始之间，观测卫星最多机动两次。同一次飞越中先后观测两颗目标时，第二颗目标之前可以不再机动。

\subsection{数据说明}

数据给出了 R1 和 30 颗目标在初始历元的位置和速度，单位为 m 和 m/s，计算时统一换算为 km 和 km/s。同步轨道半长轴由恒星日周期计算，为 $\aGEO = 42164.17$~km。由初始状态算出的轨道根数概况见表~\ref{tab:elements}，所有卫星都是同步带附近的近圆轨道。

\begin{table}[!htb]
  \centering
  \caption{初始时刻轨道根数概况}
  \label{tab:elements}
  \zihao{5}
  \begin{tabular}{lcccc}
    \toprule
     & $a-\aGEO$ / km & 偏心率 & 倾角 / ($^\circ$) & 相对 R1 的相位 / ($^\circ$) \\
    \midrule
    R1 & 1.77 & 0.00113 & 4.38 & 0 \\
    30 颗目标 & $-5.90 \sim 4.73$ & $0.00003 \sim 0.00425$ & $0.12 \sim 9.01$ & $-18.8 \sim 18.9$ \\
    \bottomrule
  \end{tabular}
\end{table}

\subsection{符号说明}

本文使用的主要符号见表~\ref{tab:symbols}，其余符号在首次出现时说明。

\begin{table}[!htb]
  \centering
  \caption{主要符号说明}
  \label{tab:symbols}
  \zihao{5}
  \begin{tabular}{lll}
    \toprule
    符号 & 含义 & 单位 \\
    \midrule
    $\mu$ & 地球引力常数 & km$^3$/s$^2$ \\
    $\aGEO$ & 同步轨道半长轴 & km \\
    $a,\ e,\ i$ & 卫星的半长轴、偏心率、倾角 & km，—，($^\circ$) \\
    $n$ & 同步轨道平均角速度 & rad/s \\
    $\Phi$ & 二体问题的状态传播映射 & — \\
    $\bm r_O,\ \bm v_O$ & 观测卫星的位置、速度矢量 & km，km/s \\
    $\bm r_b$ & 目标卫星 $b$ 的位置矢量 & km \\
    $\bm r_\odot$ & 太阳的位置矢量 & km \\
    $\bm\rho$ & 观测卫星相对目标的位置矢量，$\bm\rho = \bm r_O - \bm r_b$ & km \\
    $d$ & 观测距离，$d = \|\bm\rho\|$ & km \\
    $\alpha$ & 光照角 & ($^\circ$) \\
    $\theta$ & 侦测角 & ($^\circ$) \\
    $\beta$ & 目标处太阳方向与地心方向的夹角 & ($^\circ$) \\
    $t_k,\ \Delta\bm v_k$ & 第 $k$ 次脉冲的时刻和速度增量 & s，m/s \\
    $s$ & 观测开始时刻 & s \\
    $\tau$ & 连续观测时长，$\tau = 60$~s & s \\
    $T$ & 任务时长 & s \\
    $\lambda$ & 相对 R1 的相位 & ($^\circ$) \\
    $\Delta a$ & 半长轴相对同步轨道的偏差 & km \\
    $\delta a$ & R2 初始半长轴偏差的上限 & km \\
    $N$ & 成功观测的不同目标数 & 颗 \\
    $J$ & 累计速度增量 & m/s \\
    \bottomrule
  \end{tabular}
\end{table}
